% Complete Section 8.4, original PDF 111 / printed 101.
% block: manifold-equations-heading page:111 type:section
\subsection{{流形上的微分方程}}\label{sec:manifold-equations}

% block: finite-element-equations-on-manifolds page:111 type:ordinary-paragraph
{大部分时候, 我们取的区域是 $\mathbb R^n$ 中的有界开集. 然而, Hodge Laplace 问题可在任意 Riemannian 流形上定义, 无论有无边界, 且解在相差适当定义的调和形式的意义下唯一. 因而可以用 FEEC 求解流形上的 Hodge Laplace 问题, 但这带来大量额外考察. 流形上偏微分方程的求解本身就是重要的研究领域, 即使对 Laplace–Beltrami 方程也是如此, 它是流形上 0-形式 Laplacian 的另一名称. 一种常见的方法, 是用多面体 (即分片平坦的) 流形或更高阶分片多项式逼近来逼近流形本身, 将流形上的偏微分方程与近似流形上的偏微分方程联系起来, 再用有限元求解后者. (在剖分后的多面体流形上实现有限元方法很直接, 至少在流形嵌入 Euclidean 空间时如此.) 关于这类曲面偏微分方程的有限元, \BookCite{39} 给出了出色的综述. 对流形上 $k$-形式的 Hodge Laplacian, 当 $k>0$ 时也可采用类似方法. 然而, 误差分析要求显著推广 Chapter~\ref{chap:approximation-hilbert-complexes} 所介绍的 Hilbert 复形逼近理论. Holst 和 Stern 在 \BookCite{51} 中发展了这种推广, 以及它对 Euclidean 空间中嵌入超曲面上 Hodge Laplacian 的应用.}
